% Generated draft: requires mathematical and source review.
\documentclass[11pt]{article}
\usepackage[margin=25mm]{geometry}
\usepackage{fontspec}
\setmainfont{Latin Modern Roman}
\usepackage{amsmath,amssymb,mathtools,unicode-math}
\setmathfont{Latin Modern Math}
\usepackage{enumitem,needspace,xurl,hyperref}
\setlength{\parindent}{0pt}
\setlength{\parskip}{5pt}
\setlength{\emergencystretch}{3em}
\newcommand{\N}{\mathbb{N}}
\newcommand{\Z}{\mathbb{Z}}
\newcommand{\Q}{\mathbb{Q}}
\newcommand{\R}{\mathbb{R}}
\newcommand{\C}{\mathbb{C}}
\newcommand{\F}{\mathbb{F}}
\newcommand{\A}{\mathbb{A}}
\newcommand{\PP}{\mathbb{P}}
\newcommand{\E}{\mathbb{E}}
\newcommand{\eps}{\varepsilon}
\newcommand{\dd}{\,\mathrm{d}}
\begin{document}
\section*{UnsolvedMath ID: 20003327}
\textbf{Source problem code:} catalogue entry
\subsection{Definitions and mathematical statement}
The assigned text is a middle fragment: ``Conjecture 39 assuming the generalized Riemann hypothesis. His proof also implies the optimal rate of convergence. Unconditionally,''. It omits the conjecture, the preceding subject, and the unconditional conclusion. The aid identifies the context as arithmetic quantum unique ergodicity (QUE), Thomas C. Watson's conditional rate, and the following qualitative theorem of Lindenstrauss. The precise arithmetic surfaces and eigenfunction hypotheses are not supplied, so no theorem for arbitrary eigenfunctions or all hyperbolic surfaces is inferred.

For normalization, let \(M\) have finite hyperbolic volume, and let \(\phi_j\) be the eigenfunctions in the intended arithmetic setting with \(\int_M|\phi_j|^2\,\dd V=1\), Laplace eigenvalues \(\lambda_j\to\infty\). Put
\[
 \mu_j=|\phi_j|^2\,\dd V,\qquad m=\dd V/\operatorname{Vol}(M),\qquad
 L_j(a)=\int a\,\dd\mu_j-\int a\,\dd m.
\]
Both measures have mass one. If the eigenfunctions are instead normalized against \(m\), the defining probability must be adjusted consistently. The aid reports Watson's fixed-test estimate under the relevant generalized Lindelof-type bounds implied by GRH,
\[
 L_j(a)=O_{a,\varepsilon}(\lambda_j^{-1/4+\varepsilon}).
\]
For spectral parameter \(r_j\) with \(\lambda_j=1/4+r_j^2\), this corresponds to a power \(r_j^{-1/2+\varepsilon}\) after renaming \(\varepsilon\). The source's ``optimal'' is a supplied historical description, not a lower-bound theorem proved here. The following ``Unconditionally'' is not a claim of an unconditional optimal individual rate.

A fixed finite span does inherit a precise bound: if \(|L_j(a_i)|\leq C_iR_j\) for finitely many fixed tests \(a_1,\ldots,a_N\), then
\[
 \left|L_j\!\left(\sum_i c_ia_i\right)\right|
 \leq R_j\sum_i|c_i|C_i.
\]
This does not give a uniform bound on an entire infinite-dimensional smooth norm ball.

The aid's separate smooth-to-H\"older lemma assumes a compact smooth manifold and a uniform estimate
\[
 |L_j(a)|\leq A R_j\|a\|_{C^s}\qquad\text{for every smooth }a,
\]
with \(A\) independent of the test and \(j\), and discrepancies of probability measures. For \(0<\beta\leq1\) with \(\beta<s\), in the standard H\"older/smoothing convention, it yields
\[
 |L_j(f)|\leq C(AR_j)^{\beta/s}\|f\|_{C^\beta}.
\]
The probability bound \(|L_j(g)|\leq2\|g\|_\infty\) is part of this interpolation. The stronger uniform smooth hypothesis cannot be replaced merely by separate fixed-test big-O constants; compactness and the relevant smoothing norms are also required before applying the lemma to a QUE setting.

For a fixed test, an average variance
\[
 V_N(a)=\frac1N\sum_{j\leq N}|L_j(a)|^2
\]
only supplies the corresponding average consequences, for example
\[
 \frac1N\#\{j\leq N:|L_j(a)|>\eta\}\leq V_N(a)/\eta^2,
 \qquad\max_{j\leq N}|L_j(a)|\geq\sqrt{V_N(a)}.
\]
A positive variance asymptotic constrains such averages and maxima, but is not an individual upper rate for every eigenfunction. The exact variance normalization and asymptotic referred to by the report are not supplied. The normalization, finite-span estimate, compact interpolation lemma, and variance audit are conditional logical transfers, not a new proof of the missing Conjecture 39 or an unconditional quantitative QUE theorem.
\subsection{Short English statement}
The supplied fragment concerns a conditional arithmetic QUE rate. Its audit separates mass normalization, fixed-test constants, uniform smooth bounds, compact Hölder interpolation, and average variance information. None of these converts the unfinished unconditional sentence into an optimal rate for every eigenfunction.
\subsection{Sources}
\begin{itemize}
\item Local UnsolvedMath registry, permanent ID \(20003327\), exact source code \texttt{catalogue entry}; canonical \texttt{aim-other-notes.json}, \texttt{notes[133]}, AIM domain Other, registry category Miscellaneous, source tag \texttt{conjecture}.
\item AIM workshop ``Emerging applications of measure rigidity,'' assigned Conjecture 39 status fragment: \url{https://aimath.org/WWN/measrigid/measrigid.pdf}. The supplied verification confirms the assigned record is not a complete sentence. No PDF was fetched.
\item The aid identifies Thomas C. Watson's conditional arithmetic QUE rate and Lindenstrauss's following qualitative theorem, but supplies no full paper citations, exact arithmetic/eigenfunction hypotheses, or variance formula. The original literature field is empty.
\item Research attempt 1: result \texttt{context\_only}, problem status at run \texttt{not\_a\_problem}, rigor label \texttt{proved}, novelty confidence \texttt{low}; classification \texttt{OPEN-TRIAGE}, registry status \texttt{open}, suggested difficulty L2. Created and updated 14 August 2026. These are supplied assessments, not a quantitative resolution of an unreproduced conjecture.
\item Attribution: Ulam AI / UnsolvedMath Contributors, CC BY 4.0. The machine-generated research aid is not peer reviewed and requires independent verification before citation.
\end{itemize}
\end{document}
